\section*{第 \ref{ch:g12:brain-and-movement} 章——从意图到运动}

\begin{solution}{exo:g12:brain-and-movement:1}
每个半球顶部的一条带状区域，位于额叶与顶叶之间的沟
之前。其图谱将该区域的每一部分
对应于身体对侧的一个部位，
其面积与该部位运动的精细程度成正比。
\end{solution}

\begin{solution}{exo:g12:brain-and-movement:2}
\omterm{def:g12:brain-and-movement:cortex}{初级运动皮层}\omterm{def:g12:stretch-reflex:neuron}{神经元}的轴突穿过脑干下行，其中大多数
交叉至对侧，随后沿脊髓下行至腿部所在节段，
终止于运动\omterm{def:g12:stretch-reflex:neuron}{神经元}（或相邻的中间\omterm{def:g12:stretch-reflex:neuron}{神经元}），
其轴突延伸至腿部肌肉。
\end{solution}

\begin{solution}{exo:g12:brain-and-movement:3}
面积取决于控制的精确程度，而非大小：手能做出
许多精细且相互独立的运动，需要大量皮层\omterm{def:g12:stretch-reflex:neuron}{神经元}；
躯干做出的运动较少且较粗略。
\end{solution}

\begin{solution}{exo:g12:brain-and-movement:4}
来自皮层、\omterm{def:g12:stretch-reflex:reflex}{反射弧}和中间\omterm{def:g12:stretch-reflex:neuron}{神经元}的每一条指令
都必须经过该肌肉的运动\omterm{def:g12:stretch-reflex:neuron}{神经元}，后者汇总所有输入，
只有当总和达到阈值时才会放电；通往肌肉的
没有其他途径。
\end{solution}

\begin{solution}{exo:g12:brain-and-movement:5}
音乐家手指区域的扩大（以及志愿者经过数日练习后的扩大）；
截肢后，邻近区域接管被截肢肢体的皮层区域；
中风后，由病灶周围区域控制的恢复。
\end{solution}

\begin{solution}{exo:g12:brain-and-movement:6}
约 \qty{4}{mV}；四个相继到达的输入
可使膜电位达到 \qty{15}{mV} 的阈值；
与兴奋性输入同时到达的抑制性输入会将其抵消，电位几乎不变。
\end{solution}

\begin{solution}{exo:g12:brain-and-movement:7}
左侧的面部和手臂受损较重，腿部较轻：右侧
半球运动皮层的下部和中部（面部与手臂区域）受损，
腿部区域仅上部受到部分波及——位于交叉点以上，
因为缺损位于一侧并包括面部。
\end{solution}

\begin{solution}{exo:g12:brain-and-movement:8}
中风后\omterm{def:g12:stretch-reflex:reflex}{反射弧}本身完好，但皮层对运动\omterm{def:g12:stretch-reflex:neuron}{神经元}的抑制性
影响丧失，因此肌梭的输入不再受到抑制，
产生夸张的反应。切断神经则破坏\omterm{def:g12:stretch-reflex:reflex}{反射弧}本身：
没有信号到达肌肉，也就不会产生\omterm{def:g12:stretch-reflex:reflex}{反射}。
\end{solution}

\begin{solution}{exo:g12:brain-and-movement:9}
双腿及躯干下部的随意运动丧失（\omterm{prop:g12:brain-and-movement:pathway}{皮质脊髓束}已中断）；
腿部的\omterm{def:g12:stretch-reflex:reflex}{反射}依然存在，且更为夸张，因为其\omterm{def:g12:stretch-reflex:reflex}{反射弧}位于损伤平面以下；
面部的纤维在脑干中离开\omterm{prop:g12:brain-and-movement:pathway}{皮质脊髓束}，位于损伤平面以上，
未受影响。
\end{solution}

\begin{solution}{exo:g12:brain-and-movement:10}
紧咬牙关会使皮层向脊髓发出普遍性兴奋，其中部分
传至腿部运动\omterm{def:g12:stretch-reflex:neuron}{神经元}；它们的膜电位更接近阈值，
肌梭的输入更容易使其放电，且能使更多\omterm{def:g12:stretch-reflex:neuron}{神经元}
放电。
\end{solution}

\begin{solution}{exo:g12:brain-and-movement:11}
多年的练习通过可塑性扩大了手指区域；
在七岁之前开始练习，即皮层重组最易发生的时期，
所产生的变化比之后进行同样的练习更大。
\end{solution}

\begin{solution}{exo:g12:brain-and-movement:12}
回路因使用而得到加强：若健侧手臂承担一切活动，
通往瘫痪手臂的残存回路便从未被激活，也不会被选择；
强制使用可使其放电，而成功放电的回路将得到加强。
与小猫的情况一样，被使用的连接占据该区域，
未被使用者则失去该区域。
\end{solution}

\begin{solution}{exo:g12:brain-and-movement:13}
手部的感觉区域失去输入后，被相邻的面部区域侵占，
因此刺激脸颊会激活仍“代表”手部的\omterm{def:g12:stretch-reflex:neuron}{神经元}。
运动图谱与之相邻，也以同样方式重组：原手部区域
预计将被手臂和
面部所接管。
\end{solution}

\begin{solution}{exo:g12:brain-and-movement:14}
运动\omterm{def:g12:stretch-reflex:neuron}{神经元}受损：瘫痪，\omterm{def:g12:stretch-reflex:reflex}{反射}消失，肌肉萎缩
（失去神经输入）；皮层形成的意图无法到达
任何肌肉——可读取皮层信号的接口可驱动假肢
或肌肉刺激器。皮层受损：瘫痪，\omterm{def:g12:stretch-reflex:reflex}{反射}增强，
肌肉由\omterm{def:g12:stretch-reflex:reflex}{反射}维持；运动\omterm{def:g12:stretch-reflex:neuron}{神经元}完好，接口可
直接刺激脊髓或\omterm{def:g12:stretch-reflex:neuron}{神经元}，但意图本身已经丧失；
因此周围皮层的可塑性是主要的希望
所在。
\end{solution}

\begin{solution}{exo:g12:brain-and-movement:15}
肌肉随训练而增长，但习得的运动是神经系统的
改变：控制该运动的皮层区域扩大，其连接
变得更加精确（可塑性），到达每个运动\omterm{def:g12:stretch-reflex:neuron}{神经元}的
输入模式——哪些\omterm{def:g12:stretch-reflex:neuron}{神经元}放电、放电多少、以何种
顺序——得到优化，使运动\omterm{def:g12:stretch-reflex:neuron}{神经元}处的总和
能在正确的时刻产生正确的收缩。同样的肌肉，指挥得更好。
\end{solution}

\begin{solution}{pb:g12:brain-and-movement:1}
\textbf{1.} \omterm{prop:g12:brain-and-movement:pathway}{皮质脊髓束}纤维在脑干处交叉：左
侧皮层支配右侧。

\textbf{2.} 整条区域的下部与上部：面部和手臂（
区域侧缘）以及腿（顶部）——一处贯穿左侧
运动皮层的损伤。

\textbf{3.} 感觉区紧邻运动区，二者隔着中央沟，
未受累：感觉照常传至该区。

\textbf{4.} 脊髓中的\omterm{def:g12:stretch-reflex:reflex}{反射弧}完好，
但皮层对运动\omterm{def:g12:stretch-reflex:neuron}{神经元}的阻尼性输入已消失：
肌梭的输入更容易使\omterm{def:g12:stretch-reflex:neuron}{神经元}放电。

\textbf{5.} 右侧的运动\omterm{def:g12:stretch-reflex:neuron}{神经元}或其神经受损：\omterm{def:g12:stretch-reflex:reflex}{反射弧}
中断（无\omterm{def:g12:stretch-reflex:reflex}{反射}），肌肉失去神经输入而萎缩。
损伤位于\omterm{prop:g12:brain-and-movement:integration}{最后公路}之下，而非之上。

\textbf{6.} $12 \times 1 + 3 \times 2 = \qty{18}{mV}$：高于阈值，
故\omterm{def:g12:stretch-reflex:neuron}{神经元}放电。

\textbf{7.} $18 - 4 = \qty{14}{mV}$：刚好低于阈值，不放电
——抑制作用调节了该阶跃。

\textbf{8.} $0 + 6 = \qty{6}{mV}$：不放电；腿不运动。

\textbf{9.} 有阻尼时：$20 - 6 = \qty{14}{mV}$，低于阈值或恰在阈值处，
仅引起轻微抽动。阻尼消失后：\qty{20}{mV}，远高于阈值，引起强烈抽动：
阻尼作用正是使\omterm{def:g12:stretch-reflex:reflex}{反射}保持温和的原因。

\textbf{10.} $8 + 6 = \qty{14}{mV}$：差一点达到阈值；若
肌梭输入（或用力）再稍增，便会放电。该运动
软弱迟疑，时而放电，时而不放电。

\textbf{11.} 来自损伤周围幸存的皮层，以及来自
对侧半球的\omterm{def:g12:stretch-reflex:neuron}{神经元}，它们形成新的或加强的连接，
通向腿部运动\omterm{def:g12:stretch-reflex:neuron}{神经元}。

\textbf{12.} 可塑性会加强被激活的回路：
尝试性运动使幸存的皮层\omterm{def:g12:stretch-reflex:neuron}{神经元}及其通路放电，
产生收缩的通路被保留下来；被动移动腿部
只激活感觉回路，而不激活发出指令的
回路。

\textbf{13.} 少数未交叉的\omterm{prop:g12:brain-and-movement:pathway}{皮质脊髓束}纤维，以及
经脑干的通路，可以把指令从右侧半球
传至右腿；重组过程会募集
它们。

\textbf{14.} 手部占据一块很大的专门化区域，其精细
控制难以由邻近区域替代，且其运动
对精确度的要求远高于迈步；腿部较粗略的指令
较易重建。

\textbf{15.} 相同之处：恢复取决于使用，闲置的回路会
丧失。不同之处：成年皮层的重组远比幼猫
在关键期内的重组慢，也不够彻底——这正是
她恢复历时数月且仍不完全的原因。

\textbf{16.} 运动皮层、\omterm{prop:g12:brain-and-movement:pathway}{皮质脊髓束}、脑干中的交叉、
脊髓中的运动\omterm{def:g12:stretch-reflex:neuron}{神经元}（与\omterm{def:g12:stretch-reflex:reflex}{反射}输入的汇聚点），
神经肌肉接头处的乙酰胆碱，以及
\omterm{def:g10:muscles-and-joints:muscle}{肌纤维}\omterm{def:g10:cells-common-unit:organelle}{线粒体}产生的 \omterm{def:g12:respiration-fermentation:atp}{ATP}。

\textbf{17.} \omterm{prop:g12:glucose-and-diabetes:hormones}{胰高血糖素}；胰腺的\omterm{prop:g12:glucose-and-diabetes:hormones}{胰岛}。

\textbf{18.} 大脑的预判（运动指令的副本送至
心血管中枢，并伴随\omterm{prop:g10:heart-lungs-effort:control}{肾上腺素}），以及来自
肌肉和血液中感受器的\omterm{def:g11:hormones-and-reproduction:hormone}{反馈}。

\textbf{19.} 可塑性：参见关于视觉与大脑的章节，以及
本章；血管阻塞：参见关于心脏以及\omterm{def:g10:sport-and-health:training}{身体活动}与健康的章节；
免疫：参见关于\omterm{def:g12:innate-immunity:immunity}{先天免疫}与
\omterm{def:g12:innate-immunity:immunity}{适应性免疫}的两章。

\textbf{20.} 左侧半球的\omterm{def:g12:brain-and-movement:cortex}{初级运动皮层}；在几毫秒内
到达的输入总和使电位较静息值升高 \qty{15}{mV}；
皮层的可塑性。
\end{solution}
